\documentclass[10pt,a4paper]{amsart}
\usepackage[margin=19mm]{geometry}
\usepackage{amsmath,amssymb,mathtools}
\usepackage[scr=rsfs,cal=euler]{mathalfa}
\usepackage{xfrac}
\usepackage[unicode,hidelinks,bookmarks=false]{hyperref}
\newcommand{\ie}{i.e.\ }
\newcommand{\cf}{cf.\ }
\newcommand{\resp}{respectively\ }
\newcommand{\Subsection}{\subsection}
\providecommand{\nobreakdash}{\nobreak\hbox{-}}
\setlength{\emergencystretch}{2em}
\multlinegap0pt
\newcommand{\C}{\mathbb{C}}
\newcommand{\Z}{\mathbb{Z}}
\newcommand{\del}{\partial}
\newtheorem{theorem}{Theorem}
\newtheorem{lemma}{Lemma}
\title[Theorem 3.1 of arXiv:2508.21691]{Breaking Universality in the Lower Order Terms in the 1-level and 2-level Density of Holomorphic Cusp Newforms: the asymptotic formula for $A_{k,N}(\phi)$ (Theorem 3.1)}
\author{Lawrence Dillon, Xiaoyao Huang, Say-Yeon Kwon, Meiling Laurence, Steven J. Miller, Vishal Muthuvel, Luke Rowen, Pramana Saldin, and Steven Zanetti}
\date{Source: arXiv:2508.21691v1 (CC BY 4.0)}
\begin{document}
\maketitle
\noindent\emph{Source and licence.} Breaking Universality in the Lower Order Terms in the 1-level and 2-level Density of Holomorphic Cusp Newforms. Version arXiv:2508.21691v1 (CC BY 4.0), distributed under the Creative Commons Attribution 4.0 International licence, http://creativecommons.org/licenses/by/4.0/, as recorded in the arXiv OAI licence metadata for this exact version and shown on the abstract page https://arxiv.org/abs/2508.21691v1. Full rights notice: licenses/arxiv-2508.21691-CC-BY-4.0.txt.\par\smallskip
\noindent\textit{Editorial scope.} Counted unit: Theorem 3.1 of the paper (native source0.tex lines 471--475), the asymptotic evaluation of $A_{k,N}(\phi)$ with error $O(1/\log^4(R))$, delivered together with the paper's own complete proof of it (source0.tex lines 476--540) as the single primary proof body. No mathematics is written, repaired or re-derived: statement and proof are the paper's own text line for line, in the paper's own notation. The proof span starts at the paper's run-in proof heading \textit{Proof of Theorem 3.1.} and ends at the author's own \textbackslash qed; inside it the paper states and proves one auxiliary lemma in full (in a proof environment, retained verbatim) and states a second one for the mirrored point $z_0=k/4+1/2$, disposed of by the authors' phrase ``By the same argument, but for ... we get the following lemma'', which is retained verbatim as well. Required context retained uncounted: lines 302--305, the paper's explicit formula and, in particular, the definition $A_{k,N}(\phi):=2\hat\phi(0)\log(\sqrt{N}/\pi)+\int\psi(k/4+2\pi ix/\log(R))\phi(x)dx+\int\psi(k/4+1/2+2\pi ix/\log(R))\phi(x)dx$ that the counted theorem evaluates; the proof's final step uses exactly this definition. Omitted from this package and not redistributed: the paper's title page, abstract, subject classification and the rest of the source (all lines other than the three spans above), the paper's own preamble and title formatting, and the paper's complete bibliography (only the three entries cited inside the retained spans are transcribed). The wrapper re-states, verbatim from the paper's own preamble, the notation macros used by the retained spans ($\C$, $\Z$ and $\del$, which the paper introduces through its own \textbackslash nc alias of \textbackslash newcommand) and nothing else; the wrapper supplies its own theorem numbering, and the paper's printed identifiers are named in this scope and in the item metadata rather than silently substituted. No retained span contains a graphics-inclusion command and the paper has no figures at all, so no image asset is shipped or declared.
\section*{Retained context: the explicit formula and the definition of $A_{k,N}(\phi)$ (paper \S2)}
Our main tool is the explicit formula (see \cite{ILS2} (4.11), for example), \begin{equation}\label{eq:1lvl density breakdown}
    D_1(f, \phi) \ = \ \frac{A_{k,N}(\phi)}{\log(R)} - 2\sum_{p}\sum_{m=1}^\infty \frac{\alpha_f(p)^{m}+\beta_f(p)^{m}}{p^{m/2}} \frac{\log(p)}{\log(R) } \hat\phi\left(m \frac{\log(p)}{\log(R) }\right)
\end{equation} where 
    $$A_{k,N} (\phi)\ :=\  2\hat\phi(0)\log\left(\frac{\sqrt{N}}{\pi} \right) + \int_{-\infty}^\infty \psi\left(\frac{k}{4}+ \frac{2\pi i x}{\log(R) }\right)\phi(x)dx + \int_{-\infty}^\infty \psi\left(\frac{k}{4}+ \frac{1}{2} +\frac{2\pi i x}{\log(R)}\right)\phi(x)dx . $$

\section*{Counted claim: Theorem 3.1 of the paper, the asymptotic formula for $A_{k,N}(\phi)$, and the author's complete proof of it}
\begin{theorem}\label{gammaf}  Let $\phi$ be an even Schwartz function whose Fourier transform has compact support and let $\psi$ be the digamma function. We have the following estimate.
    \begin{align}
A_{k,N} (\phi)& \ = \ \hat\phi(0)\log(N) +\hat\phi(0)  \left(\psi\left(\frac{k}{4}\right)+\psi\left(\frac{k}{4}+ \frac{1}{2}\right)-2\log(\pi)\right) \notag \\ &~~~~- \frac{2\pi^2}{\log^2(R)}\left(\int_{-\infty}^\infty \phi(x)x^2 dx\right)\left(\psi''\left(\frac{k}{4}\right)+\psi''\left(\frac{k}{4}+ \frac{1}{2}\right)\right) +O\left(\frac{1}{\log^4(R)}\right).
\end{align} 
\end{theorem}

\textit{Proof of Theorem \ref{gammaf}.} We start by proving a lemma.
\begin{lemma}  Let $\phi$ be an even Schwartz function whose Fourier transform has compact support and let $\psi$ be the digamma function. We have the following estimate.
    \begin{equation}
        \int_{-\infty}^\infty \phi(x)\psi\left(\frac{k}{4}+ \frac{2\pi i x}{\log(R)}\right) dx \ =\ \psi\left(\frac{k}{4}\right)\hat\phi(0) -\frac{\psi^{''}(\frac{k}{4})2\pi^2}{\log^2(R)}\int_{-\infty}^\infty \phi(x)x^2 dx + O\left(\frac{1}{\log^4(R)}\right)
    \end{equation}
\end{lemma}
\begin{proof}
Using (8.363.3) of \cite{GradshteynRyzhik1965} gives us that the leading term is $\psi(\frac{k}{4})\hat\phi(0)$. Therefore, we subtract the leading term and compute the error term more accurately. 
Using the fact that $\phi$ is even
\begin{align}
    &\int_{-\infty}^\infty \phi(x)\psi\left(\frac{k}{4}+ \frac{2\pi i x}{\log(R)}\right) dx - \psi\left(\frac{k}{4}\right)\hat\phi(0)\notag\\ & \ = \
    \frac{1}{2}\int_{-\infty}^\infty \phi(x)\left[\psi\left(\frac{k}{4}+ \frac{2\pi i x}{\log(R)}\right) + \psi\left(\frac{k}{4}- \frac{2\pi i x}{\log(R)}\right) - 2\psi\left(\frac{k}{4}\right)\right]dx. \label{eq:goal 1 what we want di gamma expand}
\end{align}

Suppose $f:\Omega\to \C,$ where $\Omega$ is an open subset of $\C$. Suppose further that $z_0 \in \Omega$. Taylor expanding about the $z_0 \in \Omega$ yields
\begin{equation}\label{eq:symm Taylor expansion up to x^4 term}
    f(z_0+z) + f(z_0-z) - 2f(z_0) \ =\  f''(z_0)z^2 + O(z^4)
\end{equation} for all $z$ such that $|z|<R$ where $R$ is the radius of convergence. Moreover, we know from Cauchy's Inequality that $R \geq \text{dist}(z_0,\del\Omega)$. 
We apply the formula in \eqref{eq:symm Taylor expansion up to x^4 term} with $\Omega = \{z\in\C : \text{Re}(z) > 0\}$, $f(z) = \psi(z)$, $z_0 = \frac{k}{4}$, and $z = 2\pi ix/\log(R)$. We remark that $\text{dist}(\frac{k}{4},\del(\{z\in\C : \text{Re}(z) > 0\})) = \frac{k}{4}$ and hence $R\geq \frac{k}{4}$. Then, we get
\begin{equation}\label{eq:digamma taylor expand}
    \psi\left(\frac{k}{4} + z\right) +  \psi\left(\frac{k}{4} - z\right) - 2 \psi\left(\frac{k}{4}\right) \ =\  \psi^{''}\left(\frac{k}{4}\right) z^2 + O(z^4)
\end{equation} for all $x$ such that $|x|<k\log(R)/(8\pi)$.

Now we integrate \eqref{eq:digamma taylor expand} against the test function $\phi$ for $|x| < k\log(R)/(8\pi)$ and \eqref{eq:goal 1 what we want di gamma expand}. Then
\begin{align}\label{eq: digamma substution}
&\int_{-k\log(R)/(8\pi)}^{k\log(R)/(8\pi)}\phi(x)\psi\left(\frac{k}{4}+ \frac{2\pi i x}{\log(R)}\right) dx - \psi\left(\frac{k}{4}\right)\hat\phi(0) \notag \\ & \ = \ \frac{1}{2}\int_{-k\log(R)/(8\pi)}^{k\log(R)/(8\pi)} \phi(x)\left[-\psi''\left(\frac{k}{4}\right)\left(\frac{2\pi x}{\log(R)}\right)^2 + O\left(\frac{2\pi x}{\log^4(R)}\right)\right]dx \notag \\
& \ = \ \frac{1}{2}\int_{-k\log(R)/(8\pi)}^{k\log(R)/(8\pi)} \phi(x)\left[-\psi''\left(\frac{k}{4}\right)\left(\frac{2\pi x}{\log(R)}\right)^2\right]dx + O\left(\frac{1}{\log^4(R)}\right)\left(\int_{-k\log(R)/(8\pi)}^{k\log(R)/(8\pi)} \phi(x)x^4dx\right) \notag \\ 
& \ = \ \frac{1}{2}\int_{-k\log(R)/(8\pi)}^{k\log(R)/(8\pi)} \phi(x)\left[-\psi''\left(\frac{k}{4}\right)\left(\frac{2\pi x}{\log(R)}\right)^2\right]dx + O\left(\frac{1}{\log^4(R)}\right).
\end{align}

We now show the tail part of the integral is negligible. Consider the series expansion of $\psi(z)$ from \cite{abramowitz1972handbook}, which holds for all $z \notin \Z^{\leq 0}$:
\begin{equation}
\psi(z) \ = \ -\gamma + \sum_{n=0}^\infty \frac{z-1}{(n+1)(n+z)}
\end{equation}
where $\gamma$ is the Euler–Mascheroni constant. For $z = \frac{k}{4}+2\pi i x/ \log(R),$ we then have that for all $n \geq 1,$:
\begin{align}\label{ineq:helping convergence of tail}
  \left|  \frac{\left(\frac{k}{4}+2\pi i x/ \log(R)\right)-1}{(n+1)(n+\left(\frac{k}{4}+2\pi i x/ \log(R)\right))}\right| & \ = \ \left|\frac{8 i \pi x + (-4 + k) \log(R)}{(1 + n)\left(8 i \pi x + (k + 4n) \log(R)\right)}\right| \notag \\
  & \ \leq \ \frac{8\pi x + |k-4| \log(R)}{(1 + n)\left((k + 4n) \log(R)\right)} \notag \\
  & \ \leq \ \frac{8\pi x+|k-4|\log(R)}{n^2}.
\end{align} 
We also know that because $\phi(x)$ is Schwartz, $\int_{k\log(R)/(8\pi)}^\infty \phi(x) P(x)dx = O\left(\frac{1}{\log^B R}\right)$ for any polynomial $P(x)$. 
Therefore, using this fact along with \eqref{ineq:helping convergence of tail}, we find 
\begin{align}\label{eq: remainder part 1/R} &\left|\int_{k\log(R)/(8\pi)}^\infty  \phi(x)\left(-\gamma + \sum_{n=0}^\infty \frac{z-1}{(n+1)(n+z)}\right) dx \right|\notag\\
   & \ \leq \ (-\gamma +1)\left|\int_{k\log(R)/(8\pi)}^\infty  \phi(x)dx \right|+  \sum_{n=1}^\infty\left|\int_{k\log(R)/(8\pi)}^\infty\phi(x)\left(\frac{8\pi x+|k-4|\log(R)}{n^2}\right)dx \right| \notag\\
   & \ = \ \left(-\gamma +1+ \left(\sum_{n=1}^\infty \frac{|k-4|\log(R)}{n^2}\right)\right)\left|\int_{k\log(R)/(8\pi)}^\infty  \phi(x)dx\right| +  8\pi \left(\sum_{n=1}^\infty \frac{1}{n^2}\right)\left|\int_{k\log(R)/(8\pi)}^\infty  x\phi(x)dx\right|
   \notag \\
   & \ = \ O\left(\frac{1}{\log^4(R)}\right).
\end{align} Similarly, we know that $$\left|\int_{k\log(R)/(8\pi)}^\infty  \phi(x)\left(-\gamma + \sum_{n=0}^\infty \frac{z-1}{(n+1)(n+z)}\right) dx \right| \ = \ O\left(\frac{1}{\log^4(R)}\right).$$ 
In addition, we know that \begin{equation} \label{eq: other piece of remainder}
    \int_{k\log(R)/(8\pi)}^{\infty} \phi(x)\left[-\psi''\left(\frac{k}{4}\right)\left(\frac{2\pi x}{\log(R)}\right)^2\right]dx + \int_{-\infty}^{-k\log(R)/(8\pi)} \phi(x)\left[-\psi''\left(\frac{k}{4}\right)\left(\frac{2\pi x}{\log(R)}\right)^2\right]dx
\end{equation} is  $O\left({1}/{\log^4(R)}\right)$ because $\phi$ is Schwartz.
Therefore, combining \eqref{eq: digamma substution},  \eqref{eq: remainder part 1/R}, and \eqref{eq: other piece of remainder}, and noticing that \begin{align}
    \int_{-\infty}^{\infty} \phi(x)x^2dx = -\hat\phi''(0)
\end{align} we have that
\begin{equation}
\int_{-\infty}^\infty \phi(x)\psi\left(\frac{k}{4}+ \frac{2\pi i x}{\log(R)}\right) dx - \psi\left(\frac{k}{4}\right)\hat\phi(0) \  =  \ \frac{2\pi^2\psi''\left(\frac{k}{4}\right)}{\log^2(R)}\hat\phi''(0) + O\left(\frac{1}{\log^4(R)}\right)    
\end{equation} as desired.
\end{proof}
By the same argument, but for $z_0 = \frac{k}{4}+ \frac{1}{2}$, we get the following lemma. 
\begin{lemma} Let $\phi$ be an even Schwartz function whose Fourier transform has compact support and let $\psi$ be the digamma function. Then
      \begin{equation}
        \int_{-\infty}^\infty \phi(x)\psi\left(\frac{k}{4}+ \frac{1}{2}+ \frac{2\pi i x}{\log(R)}\right) dx \ =\ \psi\left(\frac{k}{4} + \frac{1}{2}\right)\hat\phi(0) +\frac{2\pi^2\psi^{''}(\frac{k}{4} + \frac{1}{2})}{\log^2(R)}\hat\phi''(0) + O\left(\frac{1}{\log^4(R)}\right). 
        \end{equation}
    \end{lemma}
Using the two above lemmas, we have Theorem \ref{gammaf}. \qed

\begin{thebibliography}{99}
\bibitem[ILS00a]{ILS2} H. Iwaniec, W. Luo, and P. Sarnak. \emph{Low lying zeros of families of $L$-functions}. Publications Math\'ematiques de l'IH\'ES \textbf{91} (2000), 55--131.
\bibitem[GR65]{GradshteynRyzhik1965} I. S. Gradshteyn and I. M. Ryzhik. \emph{Table of Integrals, Series, and Products}. New York: Academic Press, 1965.
\bibitem[AS72]{abramowitz1972handbook} Milton Abramowitz and Irene A. Stegun, eds. \emph{Handbook of Mathematical Functions with Formulas, Graphs, and Mathematical Tables}, 10th ed. See Chapter 6.3, pp. 258--259. New York: Dover, 1972.
\end{thebibliography}
\end{document}
